\scanpage{286}
\begin{center}- 103 -\end{center}\hfill IX

canonique de $(\check M'\otimes_{\mathbb Z}\check M)\otimes_{\mathbb Z}\mathbb Z_\ell$, d'ailleurs invariant par $\pi_o=\pi/I=\operatorname{Gal}(\overline k/k)$, puisque par définition il est invariant par transport de structure :
\[
(9.2.9)\qquad u\in(\check M'\otimes_{\mathbb Z}\check M)\otimes_{\mathbb Z}\mathbb Z_\ell)^I
\simeq\operatorname{Hom}_{\mathbb Z}(\underline M\otimes\underline M';\mathbb Z_\ell)
\simeq\operatorname{Hom}_{\mathbb Z_\ell}(\underline M_\ell\otimes\underline M'_\ell,\mathbb Z_\ell)\quad.
\]
C'est (à la notation $u$ au lieu de $u_\ell$ près) l'élément (9.1.2) que nous voulions définir.

Pour définir (9.1.2) dans le cas général, nous aurons besoin de quelques préliminaires, que nous allons développer dans les deux sections suivantes .

9.3. \uline{Extensions panachées}.

Soient $\underline C$ une catégorie abélienne, et
\[
(9.3.1)\qquad P,\ R,\ Q
\]
trois objets de $\underline C$. On s'intéresse à la classification des objets $X$ de $\underline C$, munie d'une filtration en trois crans
\[
(*)\qquad X^o=X\supset X^1\supset X^2\supset X^3=0\quad,
\]
avec des isomorphismes donnés
\[
X^o/X^1\overset{\sim}{\longrightarrow}P\quad,\quad
X^1/X^2\overset{\sim}{\longrightarrow}R,\quad
X^2/X^3=X^2\overset{\sim}{\longrightarrow}Q\quad.
\]
Posant
\[
E=X^o/X^2\quad,\quad F=X^1/X^3=X^1\quad,
\]
on voit que $X$ donne naissance à deux extensions ordinaires

(9.3.2)
\diagram{\begin{tikzcd}[column sep=small]0\arrow[r]&P\arrow[r]&F\arrow[r]&R\arrow[r]&0\end{tikzcd}}
\quad,\quad
\diagram{\begin{tikzcd}[column sep=small]0\arrow[r]&R\arrow[r]&E\arrow[r]&Q\arrow[r]&0\end{tikzcd}}
\quad,

de $R$ par $P$ et de $Q$ par $R$ respectivement. Notre point de vue sera ici

\margin{421}
\margin{421}
\margin{286}
